L'any 2006, l'exespia rus del KGB Aleksandr Litvinenko va ser víctima d'un enverinament amb poloni 210 i es va convertir en la primera víctima confirmada que moria per la síndrome de radiació aguda.

El poloni 210 és un emissor de partícules $\alpha$ que es troba a la natura i que també es pot obtenir en laboratoris nuclears.

\figura[0.25]{fig1.png}
{\centering Aleksandr Litvinenko\par}

\begin{apartat}{1}
Escriviu la reacció de desintegració del poloni 210, si sabem que en desintegrar-se produeix un isòtop del plom.
\end{apartat}

\begin{apartat}{1}
El període de semidesintegració efectiu en el cos humà del poloni 210 és de 37 dies. Si suposem que la dosi que van subministrar a Litvinenko va ser de 5~mg, quina quantitat de poloni 210 hi havia en el seu organisme quan va morir, vint dies després de l'enverinament?
\end{apartat}

\blocetiquetat{Dada:}{Símbols químics i nombres atòmics del poloni $Z(\mathrm{Po}) = 84$ i del plom $Z(\mathrm{Pb}) = 82$}
